\documentclass[12pt]{article}
\usepackage[utf8]{inputenc}
\usepackage{amsmath, amssymb, graphicx, xcolor, nicefrac}
\usepackage{geometry}
\geometry{margin=1in}

\title{The Dual Nature of Variance in Fitness: Reinterpreting Adaptive Potential as Evolutionary Constraint in Wild Populations}
\author{Walid Mawass}
\date{}

\begin{document}

\section{Methods}

\subsection{Mutation-selection-drift balance variance in log fitness}

Drift incorporates stochasticity on the frequency of a selected variant $q$. The probability that an allele with selection coefficient $s$ segregates at frequency $q$ is (cite Hernandez et al 2026)

\begin{equation}
    P(q=\frac{j}{2N} | s) = 
    \begin{cases}
        C(s)e^{4N(-s)q}q^{4N\mu_d-1}(1-q)^{4N/mu_r-1} & \text{if } j = 1, \ldots, 2N-1 \\
        e^{-2N/mu_d} \bar{t_s} & \text{if } j = 0, \ldots, 2N
    \end{cases}
\end{equation}

Assuming co-dominance, the per-site genetic variance is

\begin{equation}
    \sigma_{r_i|s}^{2} = 2(s)^2(q|s)(1-q|s)
\end{equation}

Assuming linkage equilibrium and multiplicative fitness across sites, the mean additive genetic variation summing accross $GP(s)$ sites with selection coefficient $s$ is

\begin{equation}
    E(\sigma_{r|s}^{2}) = G P(s) E(\sigma_{r_i|s}^{2}) = G P(s) \sum_{j=0}^{j=2N} \sigma_{r_i|s}^{2} P(\frac{j}{2N} | s)   
\end{equation}

Summing across all sites, the expected total genetic variance in log fitness is

\begin{equation}
    E(\sigma_{r}^{2}) = \int_{0}^{\infty} E(\sigma_{r|s}^{2}) ds  
\end{equation}


Approximating all sites in the genome with the same selection coefficient yields

\begin{equation}
    E(\sigma_{r}^{2}) = G E(\sigma_{r_i|s}^{2}) = G \sum_{j=0}^{j=2N} \sigma_{r_i|s}^{2} P(\frac{j}{2N} | s)   
\end{equation}



\end{document}